% ====================================================================
% lco_omega.tex — [LCO-Ω] LCO per-peak Ω(정규용액 커널) + 전자항 on/off 토글 (W4 초안: 최소 정확형)
%   배치: Chapter 2(LCO), \S\ref{sec:lco-hys}·\S\ref{sec:lco-electronic}·\S\ref{sec:lco-peak} 정합 신규 subsection.
%   드라이버 ch2_lco_v1.0.24. 계승 라벨: eq:lco-eqpeak·eq:lco-spinodal·eq:lco-dUdT·eq:lco-mit·eq:lco-kirchhoff·
%   eq:dSegate·sec:lco-hys-gxi·sec:lco-hys-od·sec:lco-Se·sec:lco-Se-scale.
%   신규 라벨: sec:lco-omega·eq:lco-omega-kernel·eq:lco-etoggle.
% ====================================================================
\subsection{전이별 정규용액 커널과 전자항 on/off 토글}\label{sec:lco-omega}
LCO 평형 dQ/dV(식~\eqref{eq:lco-eqpeak})는 흑연과 \emph{같은} MSMR/정규용액 커널이다. 두 정련만 남았다 --- (i)
세 feature 의 상 성격이 서로 달라 균일 폭이 아니라 \emph{전이별} $\Omega_j^\mathrm{cat}$ 로 각 peak 의 sharp/broad 를
따로 담아야 하고(per-peak $\Omega$ 커널), (ii) T1 의 전자 엔트로피는 상온 커브엔 무영향이라 on/off \emph{토글}로
분리해야 한다. 이 소절이 둘을 닫는다(#7 문구 정정 포함).

\textbf{(1) per-peak $\Omega_j^\mathrm{cat}$ 커널.} 흑연의 두-상/고용체 판정(정규용액 곡률
$\partial\mu/\partial\theta|_{1/2}=4RT-2\Omega$)은 전극 무관이므로(\S\ref{sec:lco-hys-gxi}), LCO 세 전이도 각자의
$\Omega_j^\mathrm{cat}$ 로 sharpness 가 갈린다 --- T1(MIT, 구조적 2상역)은 $\Omega_1^\mathrm{cat}>2RT$ 의
near-delta(실측 폭은 broadening 이 담음), $x\!\approx\!\tfrac12$ order--disorder(T2$\cdot$T3)는 \emph{피팅된}
$\Omega_j^\mathrm{cat}$ 이 정하는 sharp$\sim$broad 다. 균일 폭 대신 전이별 $\Omega_j^\mathrm{cat}$ 를 커널
\begin{equation}
\Big(\frac{\dd Q}{\dd V}\Big)_j^\mathrm{cat}
=\frac{Q_j^\mathrm{cat}\,F}{\big|\,RT/[\xi_j(1-\xi_j)]-2\Omega_j^\mathrm{cat}\,\big|}
\qquad(j\in\mathcal J_\mathrm{LCO})
\label{eq:lco-omega-kernel}
\end{equation}
로 쓰면(식~\eqref{eq:lco-eqpeak} 의 $\Omega$-일반화; $\Omega_j^\mathrm{cat}{=}0$ 에서 로지스틱 폴백 bit-exact), 세
feature 의 미세구조 세분이 담겨 LCO per-peak $\Omega$ 적합이 plain 대비 $+1.25$\%p 개선한다(이 세션 ablation, tier
B). 각 전이의 실제 gap 유무는 최종적으로 \emph{피팅된} $\Omega_j^\mathrm{cat}$ 이 $2RT$ 를 넘는지가 정한다
(spinodal 문턱 식~\eqref{eq:lco-spinodal}) --- 균일 폭 $w_j$ 하나로는 이 전이별 sharpness 차를 담지 못한다.

\textbf{(2) \#7 정정 --- $\Omega_j^\mathrm{cat}$ 는 상호작용 에너지, config 엔트로피는 별도 슬롯.}
\begin{keybox}
\textbf{$\Omega_j^\mathrm{cat}$ 의 의미(혼동 가드).} $\Omega_j^\mathrm{cat}$ 는 전이당 진행좌표 $\xi_j$ 의 \emph{유효
평균장 축약}(최근접 쌍 상호작용 에너지 [J/mol])이지 \emph{미시 질서구조}가 아니다. 따라서 ``$\Omega_j^\mathrm{cat}>2RT
\Leftrightarrow x{\approx}\tfrac12$ 질서상 안정''로 직결하지 않는다 --- 질서의 config 엔트로피 $\Delta S_j^\mathrm{config}$
는 평형 중심의 온도 이동(식~\eqref{eq:lco-dUdT}, $\partial U_j^\mathrm{cat}/\partial T=\Delta S_{\rxn,j}^\mathrm{cat}/F$)
에 들어가는 \emph{별도 슬롯}이고, $\Omega_j^\mathrm{cat}$ 는 gap$\cdot$sharpness 를 정하는 \emph{별도 피팅 파라미터}다.
두 몫은 서로 다른 슬롯에 존재한다(식~\eqref{eq:lco-mit} 의 config/전자 슬롯 분리와 같은 경계).
\end{keybox}
이 분리는 제일원리 클러스터 전개를 전이별 단일 평균장 쌍상호작용으로 축약할 때의 \emph{가정 차}와 같은 취지이며
(\S\ref{sec:lco-hys-od}), 각 전이의 실제 gap 유무를 미시 질서 라벨이 아니라 \emph{피팅된} $\Omega_j^\mathrm{cat}$ 로
판정한다는 점에서 config 2상역과 상호작용 에너지를 이중계산하지 않는다.

\textbf{(3) 전자항 on/off 토글.} T1 의 전자 엔트로피 $\Delta S_e(x,T)\approx-45.7$ J/(mol\,K)(MIT 골 깊이,
식~\eqref{eq:dSegate})는 \emph{상온 커브에 무영향}이다 --- 기준온도 $T_\mathrm{ref}$ 에서 동결한 상수 오프셋으로
들어가면 $U_j^\mathrm{cat}(T_\mathrm{ref})$ 가 $\Delta H$ 재기준으로 그 값을 흡수해 목표 전위 $U(298)$ 를 보존하기
때문이다(적분형 식~\eqref{eq:lco-kirchhoff}; 전자항은 $\partial U/\partial T$ 에만 작용). 곧
\begin{equation}
\boxed{\;\Delta S_e:\ \ U(298)\text{-보존 재기준으로 dQ/dV 커브 불변},\qquad
\frac{\partial U_j^\mathrm{cat}}{\partial T}\Big|_e=\frac{\Delta S_{e,j}}{F}\propto T\ \ \text{(가역열에만 작용)}\;}
\label{eq:lco-etoggle}
\end{equation}
이므로 코드는 플래그 \code{include\_electronic\_entropy}($=$\code{False} 기본)로 분리한다 --- OFF $=$ 커브 산출(전자항
잉여), ON $=$ 다온도 $\partial U/\partial T$(가역열$\cdot$회사 다온도 데이터 검증용). ``커브 산출엔 불필요$\cdot$다온도
엔트로피에서만 선택 활성''이 규약이다. 실측 정합이 이를 뒷받침한다: LCO plain MSMR $R^2{=}0.944\approx$ 흑연 $0.940$
으로, 커브 적합 품질은 전자항과 무관하다(전자항은 $T^2$ 곡률$\cdot$가역열 신호이지 상온 dQ/dV 신호가 아니다,
\S\ref{sec:lco-Se})\cite{motohashi2009}.
\begin{verifybox}
검산 --- (i) \emph{$U(298)$ 보존.} 기본 OFF 는 기존 커브 bit-exact 를 목표하나, 현 코드는 $\Delta S_e$ 상시 ON 이라
토글 OFF 가 $\Delta H$ 재기준을 요구한다(R2 게이트에서 $U(298)$ 불변 보장). (ii) \emph{부호.} 삽입 기준
$\Delta S_{e,j}<0$ 이라 $\partial U/\partial T|_e<0$(탈리튬화 시 전자 엔트로피 방출) --- 흑연의 음의 $\Delta S_\rxn$
슬롯과 같은 삽입-기준 부호. (iii) \emph{크기.} 골 깊이 $-45.7$ J/(mol\,K) 는 게이트 미분 증폭($1/\Delta x_\mathrm{MIT}
\!\approx\!20$)의 결과라 부분몰 anchor 와 척도가 다르므로 서로의 검산으로 쓰지 않는다(\S\ref{sec:lco-Se-scale} 세 양의
구분). (iv) \emph{국소성.} 게이트는 MIT 창(T1)에만 켜지고 창 밖 $\approx0$ 이라, 전자항은 T1 슬롯에만 실린다.
\end{verifybox}

% 자산(W4 R1 신규): [W4-LCOΩ-01 per-peak Ω_j^cat 커널 eq:lco-omega-kernel(+1.25%p)·#7 정정 keybox(상호작용 vs config 슬롯)]
%   [W4-LCOΩ-02 전자항 on/off 토글 eq:lco-etoggle(U(298) 보존·∂U/∂T 전용·include_electronic_entropy)·R²=0.944≈0.940]
